\chapter{Conclusões}

O presente trabalho alcançou com êxito o objetivo geral de desenvolver o DCOU, um \textit{software} \textit{web} \textit{open-source} que reúne, em uma mesma plataforma, o dimensionamento de operações unitárias em engenharia química e uma camada de recursos didáticos. Do lado do cálculo, os módulos implementados: dimensionamento de tubulações, reatores CSTR e PFR, balanços de massa, análises de escoamento e propriedades termofísicas. Foi demonstrada precisão e utilidade prática, validados utilizando referências consolidadas na literatura, sobre uma arquitetura cliente-servidor moderna e acessível via API. Do lado didático, o glossário técnico, as explicações teóricas embutidas, as visualizações interativas e os exercícios guiados consolidaram a ferramenta como apoio efetivo ao aprendizado conceitual.

A ferramenta desenvolvida preenche uma lacuna no acesso a simulações didáticas e gratuitas, especialmente em contextos acadêmicos e para profissionais autônomos, ao reunir a automação de cálculos rotineiros com recursos que promovem ativamente o aprendizado conceitual. O caráter \textit{open-source} é fundamental para sua completude futura, permitindo contribuições comunitárias que expandam o repertório de módulos e aprimorem sua robustez, tornando-o uma plataforma escalável para o ensino e a prática da engenharia química na era da Indústria 4.0.

Como perspectivas, sugere-se a adição de funcionalidades como simulações multifásicas, otimização de processos e integração com bancos de dados remotos e associando ao Padrão de Interface CAPE-OPEN, consolidando o \textit{software} como uma alternativa viável a ferramentas proprietárias.
